\section{The ratio oracle and its normalization}\label{sec:ratio}

Suppose a closed prefix $M$ has already been selected. Put
$R=\Vset\setminus M$ and retain only arcs whose two endpoints lie in $R$.
A set $S\subseteq R$ is a closure in this residual graph if and only if
$M\cup S$ is a closure in the original graph. Arcs from $R$ to $M$ already
have their prerequisites satisfied. The next subproblem is
\begin{equation}\label{eq:ratio}
 \rho^*(R)=\max\left\{\frac{p(S)}{w(S)}:
              \emptyset\ne S\subseteq R,\ S\text{ a residual closure}\right\}.
\end{equation}
We now work on one nonempty residual graph, suppress the superscript $R$,
and use $n,m$ for its size.

\subsection{Normalization removes the quotient}
The \FS{} solves the following LP:
\begin{equation}\label{eq:normalized}
\begin{aligned}
\max\quad &\sum_i p_i y_i\\
\text{s.t.}\quad &\sum_i w_i y_i=1,\\
&y_i-y_j\le0 &&(i,j)\in\Aset,\\
&y_i\ge0 &&i\in\Vset.
\end{aligned}
\end{equation}
These $y_i$ are normalized variables. They have \emph{no upper bound of one}.
For a nonempty closure $S$, the vector $\vect{y}=\ind^S/w(S)$ is feasible and has
objective $\varrho(S)$. Its coordinates can exceed one if $w(S)<1$.

\begin{theorem}[Exact ratio normalization]\label{thm:normalization}
The feasible set of \eqref{eq:normalized} is nonempty and compact. Its
optimal value is \eqref{eq:ratio}, including when this value is negative.
\end{theorem}
\begin{proof}
The full vertex set is a nonempty closure. Positivity of the weights implies
$0\le y_i\le1/w_i$, so the feasible set is bounded and closed.
For any feasible $y$, define $C_t=\{i:y_i\ge t\}$ for $t>0$.
Each $C_t$ is a closure, and the finite step function representation gives
\[
 y=\int_0^\infty\ind^{C_t}\,dt,\qquad
 1=\int_0^\infty w(C_t)\,dt,\qquad
 \vect{p}^T\vect{y}=\int_0^\infty p(C_t)\,dt.
\]
For every nonempty $C_t$, $p(C_t)\le\rho^*w(C_t)$; for an empty set both
sides vanish. Integrating proves $\vect{p}^T\vect{y}\le\rho^*$. Conversely,
$\ind^S/w(S)$ attains the ratio of any maximizing closure.
\end{proof}

\subsection{The class of linear program this is}\label{sec:class}
Three equivalent readings of \eqref{eq:normalized} fix its place among linear
programs, and each is used later.
\begin{enumerate}[label=(\alph*),leftmargin=*]
\item \textbf{One budget row over a homogeneous difference system.} The
      inequalities $y_i-y_j\le0$ are homogeneous difference constraints on a
      digraph; one dense equality row with positive coefficients is added.
\item \textbf{The base of an order cone.} The set $\{\vect{y}\ge0:y_i\le y_j\}$
      is the cone of nonnegative order-preserving vectors on the poset, whose
      extreme rays are the characteristic vectors of the nonempty closures; the
      equality slices that cone into the polytope with vertices $\ind^C/w(C)$.
      Optimizing over it therefore returns $\max_C\varrho(C)$, which is
      \cref{thm:normalization} read geometrically. It is the cone analogue of
      Stanley's order polytope, whose vertices are the characteristic vectors of
      the filters of the poset \citep{Stanley1986}.
\item \textbf{Transposed incidence plus one dense row.} Each difference
      constraint carries a $+1$ and a $-1$ in the columns of two vertex
      variables, so the matrix is the \emph{transpose} of a digraph incidence
      matrix with one dense row appended, and the right-hand side vanishes on
      the network part. This is $\mat{A}$ in \eqref{eq:standardmatrix} below.
\end{enumerate}
Reading (c) is worth keeping in mind. In the side-constrained network flow
literature the variables are arcs and the rows are vertices, and here it is the
other way round; the two formulations are linear programming duals of one
another, so the algebra transports between them, but pricing on one side
corresponds to the ratio test on the other. The same duality places
\eqref{eq:normalized} in the structural family of isotonic regression and of
total variation on a graph, which are difference systems on vertex variables as
well --- which is why \citet{Yang2026} is close prior art.

An inequality $\vect{w}^T\vect{y}\le1$ would allow $y=0$ and return zero when every
nonempty closure has negative profit. The equality in
\eqref{eq:normalized} is therefore deliberate. The outer capacitated LP
may choose nothing, but the ratio oracle must still identify the best
\emph{nonempty} closure.

\subsection{Standard form and the ratio dual}
Introduce $s_{ij}=y_j-y_i\ge0$. The standard form is
\begin{equation}\label{eq:standard}
 \vect{w}^T\vect{y}=1,\qquad y_i-y_j+s_{ij}=0\ ((i,j)\in\Aset),
 \qquad(y,s)\ge0.
\end{equation}
There are $n+m$ variables and $m+1$ independent equations: the slack
columns contain an identity matrix, and the normalization row is independent
of the arc rows. A basis has $m+1$ variables and its complement has $n-1$.
With $\mat{D}$ denoting the arc-by-vertex matrix with row $\vect{e}_i^T-\vect{e}_j^T$, write
\begin{equation}\label{eq:standardmatrix}
 \mat{A}=\begin{pmatrix}\vect{w}^T&0\\\mat{D}&\mat{I}_m\end{pmatrix},\qquad
 \vect{b}=(1,0,\ldots,0)^T,\qquad \vect{g}=(\vect{p},0).
\end{equation}
This matrix is used for independent checking, not stored densely by the
proposed efficient implementation.

The dual is
\begin{equation}\label{eq:ratio-dual}
\begin{aligned}
\min\quad &\beta\\
\text{s.t.}\quad &w_i\beta+\divg_i(\alpha)\ge p_i &&i\in\Vset,\\
&\alpha_{ij}\ge0 &&(i,j)\in\Aset,\qquad \beta\in\Rset.
\end{aligned}
\end{equation}
Here $\beta$ is free because normalization is an equality. For any closure
$S$, summing the inequalities gives
\[
 p(S)\le\beta w(S)+\sum_{i\in S}\divg_i(\alpha)\le\beta w(S).
\]
The last inequality holds because no arc leaves a closure and nonnegative
flow may enter it. Thus a feasible dual vector provides a direct certificate
of an upper bound on every closure ratio.
